\documentclass[12pt,a4paper]{article}

% ---------------------------------------------------------------------------
%  Swiss Summer Student Particle Physics Program (S3P3)
%  Project Report, July - September 2026
%  Compile with:  tectonic -X compile ProjectReport.tex   (XeTeX engine)
% ---------------------------------------------------------------------------

% --- Fonts: Times New Roman for text, Times-compatible math ---------------
\usepackage{mathptmx}          % Times math (macro level, works under XeTeX)
\usepackage{fontspec}
\setmainfont{Times New Roman}  % real Times New Roman for the text

% --- Layout --------------------------------------------------------------
\usepackage[a4paper,margin=2.5cm,headheight=60pt,headsep=16pt,footskip=36pt]{geometry}
\usepackage{graphicx}
\usepackage{xcolor}
\usepackage{enumitem}
\usepackage{fancyhdr}
\usepackage{titlesec}
\usepackage{microtype}

\definecolor{s3p3blue}{RGB}{0,176,240}   % footer highlight colour (template)

% --- Identity -----------------------------------------------------------
\newcommand{\studentname}{Giacomo Moretti}
\newcommand{\supervisorname}{Marius K\"oppel}
\newcommand{\hostgroup}{Mu3e collaboration, SciFi detector group}
\newcommand{\hostinstitute}{D-PHYS, ETH Z\"urich}
\newcommand{\projecttitle}{Operating the Mu3e SciFi Detector: Control, Monitoring and Calibration}

% --- Header / footer ---------------------------------------------------
\pagestyle{fancy}
\fancyhf{}
\fancyhead[R]{\includegraphics[height=1.6cm]{s3p3-logo.png}}
\fancyfoot[L]{\footnotesize\textbf{\textcolor{s3p3blue}{2026 Swiss Summer Student Particle Physics Program}}\ \textbar\ \studentname}
\fancyfoot[R]{\thepage}
\renewcommand{\headrulewidth}{0pt}
\renewcommand{\footrulewidth}{0pt}

% --- Section styling -------------------------------------------------
\titleformat{\section}{\normalfont\large\bfseries}{\thesection}{0.6em}{}
\titleformat{\subsection}{\normalfont\normalsize\bfseries}{\thesubsection}{0.6em}{}
\titlespacing*{\section}{0pt}{1.2em}{0.5em}
\titlespacing*{\subsection}{0pt}{1.0em}{0.4em}
\setlist[itemize]{leftmargin=1.4em,topsep=2pt,itemsep=3pt}

\setlength{\parskip}{0.5em}
\setlength{\parindent}{0pt}

\newcommand{\pill}[1]{\textit{#1}\;}   % lead-in label for a paragraph

\begin{document}

% --- Title block ---------------------------------------------------
\begin{center}
	{\large\bfseries Swiss Summer Student Particle Physics Program}\\[4pt]
	{\bfseries Project Report, July -- September 2026}\\[6pt]
	{\bfseries \projecttitle}\\[4pt]
	{\normalsize \studentname \quad$\vert$\quad Supervisor: \supervisorname}\\[1pt]
	{\normalsize \hostgroup, \hostinstitute}\\[3pt]
\end{center}
\vspace{0.6em}

% ===============================================================
\section{Summary}

The Mu3e experiment at the Paul Scherrer Institute (PSI) searches for the
charged-lepton-flavour-violating decay $\mu^+ \to e^+e^-e^+$. In the Standard
Model this decay proceeds only through a neutrino loop, with a branching ratio
of order $10^{-54}$, so it is effectively background-free and any observation
would be unambiguous evidence of new physics. The current limit,
$\mathcal{B} < 10^{-12}$ (SINDRUM, 1988), leaves four orders of magnitude for
Mu3e to explore, down to a single-event sensitivity of about $10^{-16}$.
Reaching it means running a continuous $10^{8}\,\mu/\mathrm{s}$ beam while
suppressing two backgrounds: internal conversion $\mu \to eee\nu\bar\nu$,
controlled by energy and momentum resolution, and accidental combinatorial
background, which grows with rate and is controlled by timing and vertex
resolution. The scintillating-fibre (SciFi) detector provides the timing in the
central region, with a target resolution below $200\,\mathrm{ps}$.

My project was the SciFi detector's control, monitoring and calibration system,
built end to end for the first physics run at PSI in October 2026. At the start,
SciFi was operated through a single web page (\texttt{Scifi-monitor.js}) that
had grown organically: it worked, but it was hard to trust and hard to extend,
and it lagged behind the more mature pixel-tracker tooling. I replaced it with a
structured set of Python control scripts that drive the hardware and a new
custom web page that a shift operator can rely on, both running inside the MIDAS
data-acquisition framework and communicating through its online database (ODB).

The system was built in layers: installation and bring-up at PSI; the operator
interface; the power chain (low voltage, high voltage, front-end boards); a
configuration and hardware-test hub; the two SiPM calibrations (threshold and
breakdown-voltage scans); and the saving of calibration results to three
destinations -- the live ODB, a local per-board cache, and the permanent
calibration database (CDB). By the end of the project the detector -- low
voltage, high voltage, front-end boards and readout -- is installed and working
on the experiment back-end, the calibration result-saving path is fully
implemented and wired, and the shared control-and-monitoring code library now
covers SciFi as well as the pixel detector. The main item left is a full
end-to-end test on the experiment, including a real calibration run. Alongside
the software I also took part in the on-site detector work -- cabling, power
distribution, assembly, repairs and general service work across the full Mu3e
detector -- and in the collaboration's routine of weekly reports and
coordination meetings.

\clearpage

% ===============================================================
\section{Project Details}

\subsection{State-of-the-art and Objectives}

\pill{Physics context.}
Charged lepton flavour violation (cLFV) has never been observed. Lepton flavour
is violated in the neutral sector -- neutrinos oscillate -- but in the charged
sector the Standard Model forbids it to any measurable level: $\mu^+ \to
e^+e^-e^+$ is only possible through a neutrino loop and is predicted at
$\mathcal{B} \sim 10^{-54}$~\cite{tdr}. Many extensions of the Standard Model
raise this rate to within experimental reach, which makes $\mu \to eee$ one of
the ``golden'' muon channels, alongside $\mu \to e\gamma$ and $\mu$--$e$
conversion. The strongest existing bound, $\mathcal{B} < 1.0\times10^{-12}$ at
90\%~C.L., was set by SINDRUM in 1988~\cite{sindrum}; Mu3e aims to improve on it
by four orders of magnitude, to a single-event sensitivity of about $10^{-16}$,
using the most intense continuous surface-muon beam in the world at
PSI~\cite{tdr}.

\pill{Detector and data acquisition.}
Mu3e reconstructs the three-electron final state with ultra-thin high-voltage
monolithic active pixel sensors and measures time with two scintillator systems:
the SciFi detector in the central region (target resolution $<200\,\mathrm{ps}$)
and scintillating tiles in the recurl stations ($<50\,\mathrm{ps}$)%
~\cite{bravar,vigani}. The SciFi detector is arranged in two rings, upstream and
downstream, of twelve modules each; its 96 MuTRiG readout ASICs (32 channels
each, 3072 channels in total) are served by twelve front-end boards and read out
through the MIDAS framework~\cite{koppel}, whose online database holds all
configuration and slow-control data and whose web server hosts custom operator
pages. A single module key -- side, ring position and readout end -- is used
consistently across the ODB, the detector geometry and the
calibration~\cite{wang}.

\pill{Starting point and where work was needed.}
The pixel subdetector already had a mature MIDAS integration: dedicated operator
pages and a structured library of Python control scripts. SciFi integration was
less advanced. It relied on a single monitoring page that had accreted features
over time, with control logic scattered across disconnected scripts and no
interface a shift operator could depend on during a run. With the first physics
run scheduled for October 2026, SciFi needed a coherent operator interface, a
documented and reproducible calibration procedure whose results reach the CDB,
and a control code base shared with -- rather than duplicating -- the pixel
tooling.

\pill{Contribution.}
A dependable control and monitoring system is a prerequisite for taking SciFi
data at all, and unifying the SciFi and pixel software lowers the long-term cost
of operating the experiment. Both directly support Mu3e Phase~I data taking.

\subsection{Methodology}

\pill{Environment and existing inputs.}
All work targets the MIDAS DAQ: configuration and state live in the ODB,
operator interfaces are custom HTML/JavaScript pages served by the MIDAS web
server, and device logic is Python running on the front-end machine. No new
physics data was collected -- the project builds on what the DAQ already
exposes: the ODB slow-control and configuration tree, the existing pixel control
library and the diverged SciFi scripts, and the raw per-channel data produced by
MuTRiG threshold and current--voltage scans and by the hardware tests. The fixed
hardware layout (two rings, twelve modules and front-end boards, 96 ASICs, 24
high-voltage channels) is the reference against which ODB contents were checked.
Development used a local test ODB, with checks against real hardware on the
experiment back-end.

\pill{Existing methods vs.\ what had to be developed.}
The pixel subdetector's pages and control scripts were the existing method and
were reused as a template: the same web-server/ODB binding pattern and the same
``page-first'' workflow -- mock up the layout, agree the control granularity
with the SciFi experts, then write the Python backend. What did not exist and
had to be built were the entire SciFi page and its ODB bindings (the pixel page
is not reusable given the different detector topology); SciFi coverage in the
shared control library, factored around subdetector-agnostic abstractions for
the front-end board, ASIC and channel; a configuration tab that doubles as the
detector's hardware-test hub; and the two SiPM calibrations on a common
lifecycle with a working path to permanent storage.

\pill{The interface.}
The new page separates a read-only Geometry view -- a cumulative drill-down from
super-module to module to the 32-channel ASIC grid, drawn as mirrored upstream
and downstream rings -- from the editable Power, Readout and Calibration tabs.
Every editable value is written to the ODB, from where a front-end process
applies it to the hardware; safety logic stays in Python and is never duplicated
in JavaScript. No action is ever hard-blocked: controls stay live and the page
warns rather than refusing. The largest remaining hardcoded element, the 32-cell
high-voltage board-to-module map, is flagged for a later move into Python/ODB.

\pill{Configuration and hardware tests.}
The configuration tab grew beyond parameter editing into a general
hardware-test hub -- the former Readout tab folds into it entirely. It walks the
front-end board $\to$ module $\to$ chip $\to$ channel tree, writing to a
dedicated ODB subtree, and runs three tests: PLL (each board must recover the
clock from its data link -- the recovered-clock delay histogram is flat when
locked, sharply peaked when not), PRBS (link integrity) and TDC (per-channel
test pulse).

\pill{Calibration.}
Two per-board scans share one \texttt{prepare\,/\,execute\,/\,finalise}
lifecycle. The threshold scan walks the discriminator threshold on every channel
and reads the rate curve to locate the 0.5, 1.5 and 2.5 photo-electron
crossings; its output is loaded as the operating threshold. The
breakdown-voltage scan takes the derivative of each SiPM's current--voltage
curve and fits a Gaussian to the peak, whose mean gives the breakdown voltage
and hence the operating overvoltage -- a physics constant, not just a
diagnostic. Each scan runs independently per front-end board, so any subset of
the detector can be re-calibrated on its own. A run is launched from the page,
which writes the request to the ODB; the MIDAS PySequencer picks up the
corresponding script, which opens its own MIDAS connection (with an explicit
disconnect on exit) and executes the lifecycle.

\pill{Saving.}
Results go to three destinations, one for each question an operator asks: the
live ODB (\emph{``what is the calibration right now?''}, overwritten each scan),
a local on-disk cache of the most recent calibration per board (\emph{``what did
we last calibrate, for day-to-day operations?''}), and the permanent CDB,
indexed by name and run number (\emph{``what was the calibration on a given
run?''}). The on-disk format is a small ROOT file with separate trees for
metadata, raw scan, results, conditions and channel masks. This path was
reconciled with a colleague's parallel work on the same code.

\pill{Verification.}
Each branch -- the power bring-up sequences, the hardware tests, the calibration
scans -- was built and checked on its own before being merged into the combined
system, first against the local test ODB and then against real hardware on the
experiment back-end.

\pill{Hardware and on-site work.}
The software was not developed apart from the detector. It was written inside the
collaboration's on-site campaigns at PSI, where I also worked directly on the
hardware -- cabling, power distribution, assembly and repairs -- so the power and
readout chains the page has to represent are largely the ones I helped bring up
by hand, and the interface design followed from that. The project also ran on the
collaboration's normal working rhythm throughout: weekly progress reports,
subdetector and DAQ meetings, and continuous coordination with the DAQ team, who
own the shared code the SciFi tooling now builds on.

\subsection{Results}

The three objectives -- a dependable operator interface, a reproducible
calibration procedure feeding the CDB, and a control code base shared with the
pixel detector -- were all met, and everything that had to be ready for the
October run is in place.

\pill{Operator interface.}
The new page is complete and replaces \texttt{Scifi-monitor.js} in operation.
The Power, Readout and Calibration tabs are bound to the ODB and verified against
the test ODB; the read-only Geometry view is fully functional; and readout-status
monitoring (data rate per board, run state, buffer occupancy) is integrated, so
that ``power all green'' and ``data flowing'' are presented as the two distinct
conditions they are. Where the operator previously had an unlabelled static
diagram, there is now one coherent, live view of the detector.

\pill{Configuration and hardware tests.}
The configuration tab reached beyond the planned parameter editor and became the
detector's general hardware-test entry point, absorbing the former Readout tab.
The PLL, PRBS and TDC tests all run from it and return one verdict per board or
per link.

\pill{Calibration and saving.}
Both scans run per front-end board and produce the values actually loaded for
operation (discriminator threshold and SiPM overvoltage). The result-saving path
-- the one piece that strictly had to be ready for data taking -- is fully
implemented and wired to all three destinations (live ODB, local per-board cache,
permanent CDB), with the on-disk ROOT layout finalised.

\pill{Shared code and status on the experiment.}
The common control library now covers SciFi as well as the pixel detector, so the
duplication that motivated the project is gone. On the experiment back-end, low
voltage, high voltage, front-end boards and readout are up and were operated
through the new page during the installation week at PSI. The SciFi cabling and
power distribution are wired and under management, and the assembly, repair and
service work done during the on-site periods -- part of it across the full Mu3e
detector -- leaves the apparatus in a running state.

\pill{Consistency with expectations and future developments.}
The delivered scope matches what was planned with the SciFi group at the start of
the summer; the one notable deviation -- the configuration tab growing into a
full test hub -- was in the useful direction. The principal open item is a single
end-to-end workflow test on the experiment, page through to all three save
destinations, including a real calibration run rather than test-ODB checks.
Smaller follow-ups are moving the hardcoded high-voltage board map into
Python/ODB, extending the online monitoring histograms, and wiring SciFi's own
\texttt{CalibrationFactory} to a real entry point so the scripts no longer need
their simplified direct-construction path.

\subsection{Your contributions as an S3P3 student}

I worked within the \hostgroup{} at ETH Z\"urich, in coordination with the pixel
DAQ developers, who own the shared code and whose tooling served as the
reference. The control, monitoring and calibration system described in the
previous sections is my own work, built over the summer:

\begin{itemize}
	\item I designed and implemented the new custom page end to end -- the
	read-only Geometry view and the editable Power, Readout and Calibration tabs,
	with their ODB bindings and Python backend -- and took the design decisions
	that shape it, in particular keeping all safety logic in Python and having the
	page warn rather than block.
	\item I turned the configuration tab into the hardware-test hub and
	implemented the PLL, PRBS and TDC test flows.
	\item I implemented the threshold and breakdown-voltage calibrations on their
	shared lifecycle and the full result-saving path to the ODB, the local cache
	and the CDB, including the ROOT layout, reconciling it with a colleague's
	parallel branch.
	\item I extended the common control library to cover SciFi, factoring the
	front-end-board, ASIC and channel abstractions it is built on.
\end{itemize}

\pill{Hardware, service and collaboration.}
The software was one half of the work; the other was hands-on. During the on-site
periods at PSI I did SciFi cabling and power-distribution management, mechanical
assembly, and repairs, and I designed small hardware components. A large part of
my time went into general service work on the full Mu3e detector rather than
SciFi alone, including direct contributions to other subdetectors -- assembly,
repair and maintenance tasks across the apparatus. On the organisational side I
reported progress weekly, took part in the subdetector and DAQ meetings, and
coordinated continuously with the DAQ team, whose shared code base the SciFi
tooling was built to extend rather than duplicate. This close contact with the
collaboration shaped most of the design decisions in the sections above.

% ===============================================================
\begin{thebibliography}{9}

\bibitem{tdr}
Mu3e Collaboration,
``Technical design of the Phase~I Mu3e experiment'',
arXiv:2009.11690 (2021).

\bibitem{sindrum}
U.~Bellgardt \textit{et al.} (SINDRUM Collaboration),
``Search for the decay $\mu^+ \to e^+e^+e^-$'',
Nucl. Phys. B \textbf{299} (1988) 1.

\bibitem{bravar}
A.~Bravar \textit{et al.},
``Development of the Scintillating Fibre Timing Detector for the Mu3e
Experiment'',
arXiv:2208.09906 (2022).

\bibitem{vigani}
L.~Vigani \textit{et al.},
``The Mu3e detector'',
JINST \textbf{17} (2022) C05024.

\bibitem{koppel}
M.~K\"oppel (Mu3e Collaboration),
``Mu3e Integration Run 2021'',
arXiv:2203.07855 (2022).

\bibitem{wang}
Y.~Wang (Mu3e Collaboration),
``The DAQ system of the Mu3e SciFi detector'',
SPS/APS Joint Annual Meeting (2023).

\end{thebibliography}

\end{document}
